\section*{Chapitre \ref{ch:iv1:sequences-patterns-and-aptitude} --- Suites, motifs et aptitude}
\markboth{\protect\omTocMark{}{Réponses~: Suites, motifs et aptitude}}{\protect\omTocMark{}{Réponses~: Suites, motifs et aptitude}}

\begin{solution}{iq:iv1:sequences-patterns-and-aptitude:1}
Le nombre suivant est 360.

Regardons les rapports\index{suites~: rapports}~: $6/3 = 2$, $18/6 = 3$ et $72/18 = 4$. Le multiplicateur augmente d'une unité à chaque fois, donc le suivant est 5, et $72 \times 5 = 360$.

Les autres règles sont moins solides. Les différences, 3, 12 et 54, ont pour rapports 4 et $4.5$, qui ne se stabilisent sur rien. Un polynôme passant par les quatre termes conviendrait, mais seulement parce que quatre termes permettent toujours d'en trouver un, et une règle qui comporte autant de paramètres libres que de données ne dit rien.

\iqfollow{Quel est l'ordre d'attaque face à un élément de suite~?}{Essayer les différences, puis les rapports, puis deux suites alternées, puis les puissances et les nombres premiers, puis les astuces sur les chiffres. La plupart des éléments cèdent aux trois premiers en quelques secondes. Si rien n'apparaît au bout de 20 secondes, on passe à autre chose.}
\end{solution}

\begin{solution}{iq:iv1:sequences-patterns-and-aptitude:2}
Le nombre suivant est 33.

Prenons les différences\index{suites~: différences}~: 2, 4, 6, 8. Elles augmentent de 2 à chaque fois, donc la différence suivante est 10, et $23 + 10 = 33$. Sous forme close, le terme de rang $n$ vaut $n^2 - n + 3$, et une seconde différence constante signale toujours un polynôme du second degré.

Les rapports, $5/3$, $9/5$, $5/3$, $23/15$, ne suivent aucune régularité, donc la règle sur les différences n'a pas de rival sérieux. Trois nombres déterminent un polynôme du second degré, et celui-ci est ensuite confirmé par les deux termes restants. Cette preuve supplémentaire est ce qui rend une règle crédible.

\iqfollow{Et si les différences étaient 2, 4, 8, 16~?}{Alors les différences doublent, et la suite vaut $2^n$ plus une constante. Les secondes différences permettent de distinguer les deux cas~: 2, 2, 2 contre 2, 4, 8.}
\end{solution}

\begin{solution}{iq:iv1:sequences-patterns-and-aptitude:3}
Le nombre suivant est 3125.

Les termes sont $1^1$, $2^2$, $3^3$ et $4^4$, donc le suivant est $5^5 = 3125$. L'indice tient à ce que 27 et 256 sont tous deux des puissances. Les rapports comme les différences croissent trop vite pour être utiles~: les rapports valent 4, $6.75$ et environ $9.5$, sans régularité.

Comparons d'autres familles de puissances avant de se décider. Les carrés donnent 1, 4, 9~; les cubes donnent 1, 8, 27~; et seule la règle à exposant croissant intègre aussi 256. Si 4 et 256 sont tous deux des puissances de 2, aucune règle fondée sur les puissances de 2 ne peut inclure 27.

\iqfollow{Et 1, 2, 9, 64, 625~?}{Ce sont les $n^{n-1}$~: $1^0$, $2^1$, $3^2$, $4^3$, $5^4$. Le suivant est $6^5 = 7776$.}
\end{solution}

\begin{solution}{iq:iv1:sequences-patterns-and-aptitude:4}
La lettre suivante est N.

Ce sont les initiales des mots One, Two, Three, Four, Five, Six, Seven et Eight, donc la suivante est l'initiale de Nine.

Lues comme des rangs dans l'alphabet, les lettres donnent 15, 20, 20, 6, 6, 19, 19, 5, ce qui ne correspond à aucune règle arithmétique. Lorsqu'une suite de lettres refuse de se comporter comme des nombres et répète les lettres de façon irrégulière, elle est en général tirée d'une liste que tout le monde connaît~: noms de nombres, mois, jours de la semaine ou planètes. Le critère est que la règle explique chaque terme et ne comporte aucun paramètre libre.

\iqfollow{Quelle lettre vient ensuite dans M, T, W, T, F, S~?}{S, pour Sunday~: les jours de la semaine, à partir du lundi.}
\end{solution}

\begin{solution}{iq:iv1:sequences-patterns-and-aptitude:5}
Le terme suivant est U21.

Les nombres sont 1, 3, 6, 10 et 15, les nombres triangulaires\index{nombres triangulaires}, avec pour différences 2, 3, 4, 5. Le suivant est donc 21. Les lettres suivent la même règle, puisque chaque lettre est simplement celle qui occupe ce rang dans l'alphabet~: A est la première lettre, C la troisième, F la sixième, J la dixième, O la quinzième, et U la vingt et unième.

Repérer que la lettre et le nombre expriment le même fait est l'intérêt de l'élément. Traitées comme deux suites distinctes, elles demandent deux fois plus de travail. Aucune règle séparée pour les lettres ne serait plus simple.

\iqfollow{Quel terme vient après U21~?}{La différence suivante est 7, ce qui donne 28. L'alphabet ne compte que 26 lettres, donc l'élément doit soit s'arrêter, soit reboucler sur B. C'est pourquoi l'élément s'arrête là où il s'arrête.}
\end{solution}

\begin{solution}{iq:iv1:sequences-patterns-and-aptitude:6}
Le nombre suivant est 16.

Il s'agit de deux suites prises en alternance\index{suites~: entrelacées}. Les positions impaires donnent 1, 2, 4, 8, qui doublent. Les positions paires donnent 20, 17, 14, 11, qui diminuent de 3. Le terme suivant est en position impaire, donc il vaut $8 \times 2 = 16$.

L'indice tient à ce que la suite monte et descend en alternance, alors que les différences entre voisins (19, $-18$, 15, \dots) ne suivent aucune régularité. La lecture des termes en alternance donne deux règles très simples, chacune confirmée sur trois étapes.

\iqwrong{Répondre 8, qui est le terme suivant en position paire, $11 - 3$. Il faut compter la position du terme suivant~: la suite comporte huit termes jusqu'ici, donc le neuvième se trouve dans la suite impaire.}
\iqfollow{S'il y a trois fils au lieu de deux, comment le remarqueriez-vous~?}{Des différences espacées de trois paraissent régulières alors que l'espacement de deux ne l'est pas. Essayer d'abord un espacement de deux, puis de trois.}
\end{solution}

\begin{solution}{iq:iv1:sequences-patterns-and-aptitude:7}
Le nombre suivant est 20.

Les pas alternent~: multiplier par 2, puis retrancher 2. Donc $3 \to 6 \to 4 \to 8 \to 6 \to 12 \to 10$, et le pas suivant est une multiplication par 2, ce qui donne 20.

On peut aussi lire cette suite comme deux fils, 3, 4, 6, 10 et 6, 8, 12. Le premier a pour différences 1, 2, 4 et le second 2, 4, donc les deux différences doublent, et le terme suivant du fil est $12 + 8 = 20$. Les deux lectures donnent 20, parce qu'elles décrivent la même suite. La règle alternée est plus simple~: elle ne comporte que deux pas, alors que les fils exigent deux valeurs de départ et un doublement.

\iqfollow{Et le terme d'après~?}{Retrancher 2, donc 18. Dans la lecture en fils, il vaut $10 + 8 = 18$, et les deux points de vue concordent toujours.}
\end{solution}

\begin{solution}{iq:iv1:sequences-patterns-and-aptitude:8}
La figure manquante est un cercle gris.

Chaque ligne et chaque colonne contient une forme de chaque sorte et une couleur de chaque sorte, ce qui est la règle derrière de nombreux éléments matriciels~: un carré latin\index{carré latin}. La ligne 3 contient déjà un carré et un triangle, elle a donc besoin d'un cercle. Elle contient déjà le noir et le blanc, elle a donc besoin du gris. La troisième colonne le confirme~: elle contient un triangle et un carré, donc elle a besoin d'un cercle, et le noir et le blanc, donc elle a besoin du gris.

\iqwrong{Traiter les lignes comme un motif à recopier d'une ligne à l'autre, par exemple en décalant la ligne 1 d'une case pour obtenir la ligne 2. Cette lecture convient aux formes, mais les couleurs ne se déplacent pas de la même façon, si bien qu'il faut lire séparément les deux caractéristiques.}
\iqfollow{Quelle est la méthode pour une grille plus difficile~?}{Séparer les caractéristiques~: forme, couleur, nombre et position. Résoudre chacune isolément. Une grille n'est difficile que si l'on essaie de voir toutes les caractéristiques à la fois.}
\end{solution}

\begin{solution}{iq:iv1:sequences-patterns-and-aptitude:9}
Non, la conclusion ne découle pas des prémisses.

Il faut chercher une situation où les prémisses sont vraies et la conclusion fausse. Supposons que l'entreprise ait un seul analyste, Ann, et un seul trader, Tom, et que tous deux aient un diplôme. Alors tous les analystes ont un diplôme, et certaines personnes diplômées sont des traders, pourtant aucun analyste n'est trader. Un seul contre-exemple tranche la question\index{contre-exemple}.

Le piège est que la conclusion peut très bien être vraie dans une vraie entreprise. L'énoncé demande seulement si elle doit l'être.

\iqfollow{Remplacez la seconde prémisse par «~certains analystes sont traders~». Que s'ensuit-il~?}{Alors certaines personnes diplômées sont traders, puisque ces analystes ont un diplôme. Cette conclusion découle bien, dans tous les cas.}
\iqfollow{Remplacez-la par «~aucune personne diplômée n'est trader~». Que s'ensuit-il~?}{Alors aucun analyste n'est trader, puisque tout analyste a un diplôme.}
\iqwhy{Les sections de raisonnement testent cela, et la relecture d'une affirmation de recherche aussi~: «~les gagnants ont la caractéristique $X$, et certaines actions ayant la caractéristique $X$ sont bon marché~» ne signifie pas que les gagnants sont bon marché.}
\end{solution}

\begin{solution}{iq:iv1:sequences-patterns-and-aptitude:10}
Seule la quatrième affirmation est vraie.

Les cinq affirmations s'excluent mutuellement, donc au plus une peut être vraie. Si aucune ne l'était, les cinq seraient fausses, et la cinquième serait alors vraie, ce qui est contradictoire. Il y en a donc exactement une de vraie, et les quatre autres sont fausses. Exactement quatre fausses est ce que dit la quatrième affirmation, c'est donc elle qui est vraie.

Toute autre attribution échoue. Si seule la deuxième était vraie, il y aurait quatre affirmations fausses, et non deux.

\iqwrong{Répondre que la cinquième est vraie parce qu'elles semblent toutes fausses. Si la cinquième était vraie, les cinq seraient fausses, y compris elle-même, elle ne peut donc pas être vraie.}
\iqfollow{Et si la liste comporte $n$ affirmations de ce type~?}{Le même raisonnement fonctionne~: une seule est vraie, $n - 1$ sont fausses, et l'affirmation numéro $n - 1$ est la vraie.}
\end{solution}

\begin{solution}{iq:iv1:sequences-patterns-and-aptitude:11}
Il transforme 5 en 126.

La règle est $n^3 + 1$~: $1 + 1 = 2$, $8 + 1 = 9$, $27 + 1 = 28$, $64 + 1 = 65$, et donc $125 + 1 = 126$. L'indice tient à ce que 9, 28 et 65 valent chacun un cube plus un, et ces nombres sont assez proches pour être reconnus.

Les différences le confirment~: 7, 19, 37, puis 12, 18, puis une constante 6, ce que donne un polynôme de degré trois. En prolongeant ce tableau, on obtient 24, puis 61, et $65 + 61 = 126$. Le tableau des différences conduit donc à la même réponse sans qu'il soit nécessaire de repérer les cubes. Une règle de la forme «~multiplier par un coefficient, puis ajouter une constante~», $an + b$, échoue immédiatement, puisque les premières différences ne sont pas constantes.

\iqfollow{Pourquoi un tableau dont la troisième différence est constante implique-t-il un polynôme de degré trois~?}{La différenciation abaisse le degré d'un polynôme d'une unité, donc l'opération inverse l'élève d'une unité. On part de la constante et on la somme trois fois, ce qui donne un polynôme de degré trois. Pour $n^3$, la constante vaut $3! = 6$.}
\end{solution}

\begin{solution}{iq:iv1:sequences-patterns-and-aptitude:12}
FOX s'écrit 45.

Additionner les rangs des lettres dans l'alphabet. Pour CAT, $3 + 1 + 20 = 24$. Pour DOG, $4 + 15 + 7 = 26$. Donc FOX vaut $6 + 15 + 24 = 45$.

Les autres codes naturels échouent. Les produits donnent $3 \times 1 \times 20 = 60$ pour CAT, et compter depuis la fin de l'alphabet donne $24 + 26 + 7 = 57$. Une seule règle simple convient aux deux mots, et deux mots suffisent à écarter une règle qui ne convient à un seul que par hasard~: la somme pour CAT seul aurait pu être une coïncidence, mais qu'on retrouve la même somme sur DOG est peu vraisemblable.

\iqfollow{45 est-il propre à FOX~?}{Non. Des mots différents peuvent partager une somme, si bien qu'un codage de ce type ne peut pas être décodé en sens inverse. TAX donne aussi $20 + 1 + 24 = 45$.}
\end{solution}

\begin{solution}{iq:iv1:sequences-patterns-and-aptitude:13}
La vitesse moyenne est de 48~km/h.

La vitesse moyenne est le quotient de la distance totale par le temps total. Le trajet aller prend $120/40 = 3$ heures, et le trajet retour $120/60 = 2$ heures, donc la voiture parcourt 240~km en 5~heures~: 48~km/h. C'est la moyenne harmonique\index{moyenne harmonique} des deux vitesses, $2/(1/40 + 1/60) = 48$. On remarque que la distance n'intervient pas, puisqu'elle se simplifie.

\iqwrong{Répondre 50, qui est la moyenne simple des vitesses. La voiture passe plus de temps à la vitesse la plus faible, la moyenne doit donc être inférieure à 50.}
\iqfollow{Vous dépensez 1000~\$ en achetant à 40, puis 1000~\$ en achetant à 60. Quel est votre prix moyen par action~?}{Le même calcul, toujours 48, puisque l'on achète $25 + 16.67$ actions avec 2000~\$. À somme égale, on achète davantage d'actions au cours le plus bas, ce qui explique que l'investissement programmé donne un prix moyen inférieur à la moyenne simple.}
\end{solution}

\begin{solution}{iq:iv1:sequences-patterns-and-aptitude:14}
Le réservoir se remplit en 4~heures.

Additionner les débits, et non les durées\index{additivité des débits}. En fractions de réservoir par heure, le tuyau A apporte $1/6$, le tuyau B apporte $1/3$ et la vidange retire $1/4$~:
\[
\frac16 + \frac13 - \frac14 = \frac{2 + 4 - 3}{12} = \frac14.
\]
À raison d'un quart de réservoir par heure, un réservoir plein se remplit en 4~heures.

\iqwrong{Faire la moyenne ou additionner les durées, par exemple $6 + 3 - 4 = 5$ heures. Un tuyau qui remplit seul en 3 heures doit, avec un autre, aller plus vite que 3 heures, et une formule qui additionne les durées ne peut pas le montrer.}
\iqfollow{Que se passerait-il si la vidange ne prenait que 2~heures~?}{Le débit serait $\tfrac16 + \tfrac13 - \tfrac12 = 0$, et le réservoir resterait vide. Une vidange plus rapide encore signifie que le niveau ne monte jamais.}
\iqfollow{Et avec les seuls tuyaux A et B, sans la vidange~?}{Le débit vaut $\tfrac12$ réservoir par heure, soit 2~heures.}
\end{solution}

\begin{solution}{iq:iv1:sequences-patterns-and-aptitude:15}
Dix étudiants suivent les deux cours. Sans le nombre d'étudiants qui n'en suivent aucun, le nombre de ceux qui les suivent tous deux pourrait être n'importe quel entier compris entre 3 et 18.

Utiliser l'inclusion--exclusion\index{inclusion--exclusion}. Comme 7 n'en suivent aucun, 33 en suivent au moins un. La somme $25 + 18 = 43$ compte deux fois ceux qui suivent les deux cours, donc $43 - 33 = 10$ les suivent tous deux.

Pour l'intervalle, le recouvrement est maximal lorsque les 18 étudiants de physique suivent aussi les mathématiques, ce qui laisse 15 étudiants qui n'en suivent aucun. Il est minimal lorsque le moins possible d'étudiants n'en suivent aucun~: les 40 en suivent au moins un, ce qui fait $25 + 18 - 40 = 3$ étudiants dans les deux cours.

\iqwrong{Répondre $25 + 18 - 40 = 3$ en oubliant les 7. C'est la borne inférieure, et non la réponse.}
\iqfollow{Avec une troisième matière, la chimie, qu'est-ce qui change~?}{On additionne les trois effectifs, on retranche les recouvrements deux à deux et on rajoute le recouvrement triple. Il faut alors davantage de données, et le même raisonnement fournit encore des bornes si une partie d'entre elles manque.}
\end{solution}

\begin{solution}{iq:iv1:sequences-patterns-and-aptitude:16}
La marge du produit Z a le plus progressé, de 6~points de pourcentage. Le bénéfice du produit Y a le plus augmenté, de 25~\$.

Calculer les deux marges de chaque produit. X passe de $50/200 = 25~\%$ à $70/250 = 28~\%$, soit 3 points de plus. Y passe de $80/400 = 20~\%$ à $105/500 = 21~\%$, soit 1 point de plus. Z passe de $24/120 = 20~\%$ à $39/150 = 26~\%$, soit 6 points de plus. En dollars, les variations de bénéfice valent $+20$ pour X, $+25$ pour Y et $+15$ pour Z.

Les deux questions n'ont pas la même réponse, ce qui est courant dans ce type de question. Une base élevée peut donner la plus forte variation en dollars en même temps que la plus faible variation de marge.

\iqwrong{Comparer les croissances du bénéfice en pourcentage, ce qui est une autre question~: celle de Z, $62.5~\%$, dépasse celle de X, $40~\%$, et celle de Y, $31~\%$. Vérifier quelle des trois mesures l'énoncé demande réellement.}
\iqfollow{Comment faites-vous pour aller vite~?}{Arrondir à des fractions simples, et ne calculer que les deux candidats qui comptent. Lors d'un test, on n'a que rarement besoin des six valeurs exactes.}
\end{solution}

\begin{solution}{iq:iv1:sequences-patterns-and-aptitude:17}
Le nombre suivant est 221.

Chaque terme est le produit de deux nombres premiers consécutifs\index{suites: nombres premiers}~: $2 \times 3 = 6$, $3 \times 5 = 15$, $5 \times 7 = 35$, $7 \times 11 = 77$, $11 \times 13 = 143$, et le terme suivant est $13 \times 17 = 221$.

Les différences sont 9, 20, 42 et 66, et ni leurs rapports ni leurs écarts ne se ramènent à quelque chose de simple. Factoriser est l'étape à tenter lorsque les différences ne donnent rien et que les nombres ne sont pas des puissances. Ici chaque terme se décompose en deux nombres premiers, et le second facteur d'un terme est le premier facteur du suivant, ce qui est un motif trop fort pour être dû au hasard.

\iqfollow{Quelle serait la règle derrière 4, 9, 25, 49, 121~?}{Les carrés des nombres premiers, donc le terme suivant est $13^2 = 169$. Ici, une simple règle «~carrés~» est trop faible, car elle prédirait 16 après 9.}
\end{solution}

\begin{solution}{iq:iv1:sequences-patterns-and-aptitude:18}
Le nombre suivant est 73.

Chaque terme est égal au terme précédent augmenté de la somme de ses propres chiffres\index{suites: règles sur les chiffres}. Ainsi $19 + 10 = 29$, $29 + 11 = 40$, $40 + 4 = 44$, $44 + 8 = 52$, $52 + 7 = 59$, et $59 + 14 = 73$.

Les écarts, 10, 11, 4, 8, 7, montent et descendent sans motif arithmétique, et c'est l'indice qui invite à regarder les chiffres~: un pas qui dépend des chiffres paraît erratique, mais il est entièrement déterminé. Aucune règle portant sur les seuls écarts n'explique ces valeurs avec moins d'outils, car une telle règle exigerait autant de nombres libres qu'il y a de pas.

\iqfollow{Et pour 12, 15, 21, 24, 30, 33~?}{Deux règles conviennent~: ajouter la somme des chiffres, ou ajouter alternativement 3 et 6. Toutes deux donnent 39 ensuite, et elles divergent au coup suivant~: 51 contre 42. Un énoncé qui s'arrête là a donc une seule réponse, mais pas une seule règle, et c'est un énoncé faible.}
\end{solution}

\begin{solution}{iq:iv1:sequences-patterns-and-aptitude:19}
Le nombre suivant est 81.

Ce sont trois suites entrelacées\index{suites: entrelacées}. Les positions 1, 4, 7, 10 portent 2, 4, 8, 16, qui double à chaque fois. Les positions 2, 5, 8 portent 3, 9, 27, qui triple. Les positions 3, 6, 9 portent 5, 7, 9, qui croît de 2 à chaque fois. Le onzième terme est dans la deuxième suite, donc il vaut $27 \times 3 = 81$.

Essayer d'abord un espacement de deux~: 2, 5, 9, 8, 9 et 3, 4, 7, 27, 16 ne donnent rien. Un espacement de trois fournit trois règles très nettes, et c'est trop d'ordre pour être dû au hasard. Un indice pour repérer les fils est que les petits et les grands nombres alternent sans cesse, et que des nombres comme 9 apparaissent deux fois dans des rôles différents.

\iqwrong{Répondre 11, qui prolonge le troisième fil. La suite comporte dix termes pour l'instant, et $11 = 3 \times 3 + 2$, donc le onzième terme appartient au deuxième fil.}
\iqfollow{Quel est le terme qui suit 81~?}{Le troisième fil passe à 11.}
\end{solution}

\begin{solution}{iq:iv1:sequences-patterns-and-aptitude:20}
Le nombre suivant est 21.

Il s'agit d'un comptage en base 6\index{bases de numération}~: aucun chiffre ne dépasse 5, et après 15 vient 20, comme après 5 vient 10. Donc après 20 vient 21.

Les chiffres manquants sont l'indice. La suite saute 6, 7, 8, 9, puis 16, 17, 18, 19, ce qui se produit si ces chiffres n'existent pas. Une règle portant sur les seuls écarts, cinq pas de 1 puis un saut de 5, doit traiter le début de la suite comme un cas particulier, et elle ne dit rien de ce qui suit 55. La base 6 prévoit chaque saut, y compris celui de 55 à 100.

\iqfollow{Quelle est la valeur du terme 20 en écriture décimale~?}{$2 \times 6 + 0 = 12$. C'est le douzième terme de la suite.}
\iqfollow{Et en base 3, que suivrait 1, 2, 10, 11, 12, 20, 21, 22~?}{100.}
\end{solution}

\begin{solution}{iq:iv1:sequences-patterns-and-aptitude:21}
La case manquante porte le haut, la droite et le bas.

Dans chaque ligne, la troisième case porte exactement les côtés qui figurent dans l'une des deux premières cases, mais pas dans les deux~: il s'agit d'un ou exclusif\index{ou exclusif}. Dans la ligne 1, le haut n'apparaît que dans la première case, et la droite seulement dans la seconde, tandis que la gauche figure dans les deux et n'est pas retenue. Dans la ligne 2, le bas n'apparaît que dans la première case et le haut seulement dans la seconde, tandis que la droite figure dans les deux. Dans la ligne 3, le haut et le bas n'apparaissent que dans la première, la droite seulement dans la seconde, et la gauche dans les deux. On obtient donc le haut, la droite et le bas.

Les deux autres règles usuelles échouent sur la ligne 1. L'union contiendrait la gauche, et l'intersection ne donnerait que la gauche.

\iqwrong{Utiliser l'union pour la ligne 3, ce qui donne les quatre côtés. Tester toujours une règle sur chaque ligne complète avant de l'appliquer à la ligne incomplète.}
\iqfollow{Est-ce la même opération que dans le code~?}{Oui. Si chaque côté est un bit d'un nombre sur quatre bits, la troisième case est le XOR bit à bit des deux premières.}
\end{solution}

\begin{solution}{iq:iv1:sequences-patterns-and-aptitude:22}
Le nombre manquant est 29.

La règle est $a^2 + b^2$~: $4 + 9 = 13$, $16 + 1 = 17$, $1 + 25 = 26$, $9 + 9 = 18$, et donc $25 + 4 = 29$.

L'indice tient à ce que les grands nombres comptent bien plus que les petits~: la ligne $(1, 5)$ donne 26, alors que $(4, 1)$ ne donne que 17. Cela plaide contre une somme ou un produit simples, et en faveur des carrés. En effet, aucune règle linéaire $pa + qb + c$ ne convient aux quatre lignes, car trois lignes fixent $p$, $q$ et $c$, et la quatrième échoue alors. Un produit plus une constante échoue aussi, puisque $2 \times 3 = 6$ et $4 \times 1 = 4$ exigeraient des constantes valant 7 et 13.

\iqwrong{Ajuster une règle linéaire à deux ou trois lignes, puis s'arrêter. Trois lignes ajustent toujours trois nombres libres, donc cela ne prouve rien. C'est la quatrième ligne qui teste la règle.}
\iqfollow{Quelle quantité de preuves suffit~?}{Une règle est crédible lorsque les lignes sont plus nombreuses que les nombres libres qu'elle contient. La règle $a^2 + b^2$ n'a aucun nombre libre, et elle convient à quatre lignes.}
\end{solution}

\begin{solution}{iq:iv1:sequences-patterns-and-aptitude:23}
Seul B dit la vérité. A et C sont menteurs.

Il n'y a que huit cas, et le statut de A les détermine tous. Supposons que A dise la vérité. Alors B ment, donc son affirmation est fausse, et C dit la vérité. Mais alors l'affirmation de C, selon laquelle A et B mentent tous deux, est vraie, ce qui contredit le fait que A dise la vérité. Donc A ment. Son affirmation est alors fausse, si bien que B dit la vérité. L'affirmation de B est donc vraie, et C ment. Enfin, l'affirmation de C, selon laquelle A et B mentent tous deux, doit être fausse, et elle l'est, puisque B dit la vérité.

\iqwrong{S'arrêter à la première attribution cohérente, sans tester l'autre valeur de A. Une énigme de ce type peut avoir deux solutions, comme le montre la relance, il faut donc tester les deux.}
\iqfollow{Et si C disait plutôt que A dit la vérité~?}{Les deux cas sont alors cohérents~: A et C disent la vérité avec B menteur, ou A et C mentent avec B qui dit la vérité. La question n'aurait pas de réponse unique.}
\end{solution}

\begin{solution}{iq:iv1:sequences-patterns-and-aptitude:24}
L'ordre est C, A, B, E, D.

Commencer par les faits les plus forts. La chaîne C devant A devant B fixe trois des cinq positions relatives les unes aux autres. Le couple E et D est soudé, E étant immédiatement devant D, on le traite donc comme un seul bloc. Le bloc a alors quatre places possibles~: avant C, entre C et A, entre A et B, ou après B. Comme B n'est pas dernier, quelque chose suit B. Seul le bloc peut le suivre, donc c'est lui qui vient en dernier.

Il existe $5! = 120$ ordres au total. La chaîne en conserve $120/6 = 20$, le bloc en conserve 4 d'entre eux, et l'indice final n'en laisse exactement qu'un.

\iqwrong{Placer E et D sans lire le mot «~immédiatement~». Dans ce cas, E devant D seul permettrait par exemple E, C, A, B, D, où B n'est pas non plus dernier.}
\iqfollow{Et si l'indice «~B ne finit pas dernier~» est retiré~?}{Les quatre places du bloc sont alors permises, donc quatre ordres subsistent, et la question n'a pas de réponse unique.}
\end{solution}

\begin{solution}{iq:iv1:sequences-patterns-and-aptitude:25}
Anna a 36 ans et Ben en a 27.

Notons $A$ et $B$ leurs âges, avec $A > B$. Il y a $A - B$ ans, Anna avait l'âge actuel de Ben. À cette époque, l'âge de Ben était $B - (A - B) = 2B - A$. L'énoncé indique que $A = 2(2B - A)$, donc $3A = 4B$. Avec $A + B = 63$, cela donne $A = 36$ et $B = 27$.

Vérification par le récit~: il y a neuf ans, Anna avait 27 ans, soit l'âge actuel de Ben, et celui-ci en avait 18~; or 36 est bien le double de 18.

\iqwrong{Lire «~était~» comme «~est~» et écrire $A = 2B$. La phrase renvoie à l'âge passé de Ben, donc une différence d'âges doit apparaître. Cette différence reste la même au fil du temps, ce qui rend ces questions faciles à mettre en équation.}
\iqfollow{Et si la somme valait 70~?}{Alors $A = 40$ et $B = 30$. Le rapport $4:3$ tient pour toute somme.}
\end{solution}

\begin{solution}{iq:iv1:sequences-patterns-and-aptitude:26}
La question n'est pas défectueuse en pratique, et il faut répondre 32. Une formule donnant 31 existe, car une formule peut être construite pour donner n'importe quoi.

Cinq points fixent un polynôme de degré quatre. Par les points 1, 2, 4, 8, 16, ce polynôme donne exactement 31 au terme suivant\index{interpolation de Lagrange}. En ajoutant $c\,(n - 1)(n - 2)(n - 3)(n - 4)(n - 5)$, qui s'annule aux cinq places données, on peut donner au sixième terme n'importe quelle valeur. Aucune question sur les suites n'a donc de réponse unique au sens strict de la logique.

Ce que demande une question, c'est la règle comportant le moins de nombres libres qui convienne. Le doublement exige une valeur de départ et un multiplicateur, et il explique cinq termes. Le polynôme utilise cinq nombres pour expliquer cinq termes~; il n'explique donc rien.

\iqfollow{Le nombre 31 apparaît-il naturellement~?}{Oui. On prend $n$ points sur un cercle, placés de sorte que jamais trois cordes ne se coupent en un même point, et on joint chaque paire~: le disque est alors découpé en 1, 2, 4, 8, 16, 31 régions, pour $n = 1$ à 6. Ce nombre vaut $\binom{n}{4} + \binom{n}{2} + 1$.}
\iqwhy{Le même raisonnement fait la différence entre un backtest et un surajustement. Un modèle ayant autant de nombres libres que de points de données s'ajuste toujours. Ce qui compte, c'est combien il explique par nombre libre.}
\end{solution}

\enlargethispage{-3\baselineskip}%
\begin{solution}{iq:iv1:sequences-patterns-and-aptitude:27}
Le nombre suivant est 9.

Ce sont les sommes des chiffres\index{suites : règles sur les chiffres} des nombres carrés~: $1$, $4$, $9$, $16 \to 7$, $25 \to 7$, $36 \to 9$, $49 \to 13$, $64 \to 10$, et ensuite $81 \to 9$.

Le début de la suite 1, 4, 9 évoque des carrés, mais 16 manque. Ensuite, 7, 7, 9, 13, 10 ne suivent aucune règle arithmétique, et la chute de 13 à 10 est l'indice~: une formule régulière ne ferait pas cela, alors qu'une opération sur les chiffres le fait. On revient donc au début~: on cherche ce qui transforme 16 en 7, et on trouve la somme des chiffres.

Les différences échouent d'emblée ici. Leurs valeurs, 3, 5, $-2$, 0, 2, 4, $-3$, changent de signe sans motif, et un polynôme qui passe par huit termes n'explique rien.

\iqfollow{Quel est le terme après 9~?}{$10^2 = 100$, donc 1. La forte baisse relève encore du même effet.}
\end{solution}

\begin{solution}{iq:iv1:sequences-patterns-and-aptitude:28}
Le nombre suivant est 3411.

Chaque terme est le produit des deux précédents, plus un\index{suites : récurrences}~: $1 \times 2 + 1 = 3$, $2 \times 3 + 1 = 7$, $3 \times 7 + 1 = 22$, $7 \times 22 + 1 = 155$. Le suivant est donc $22 \times 155 + 1 = 3410 + 1 = 3411$.

L'indice tient à la croissance. Les rapports valent 2, $1.5$, $2.33$, $3.14$, $7.05$, et une accélération aussi rapide signifie que les termes se multiplient entre eux plutôt que d'être multipliés par des constantes. Il faut donc tester un produit de deux termes précédents, avec une petite correction. Cette règle n'utilise qu'un seul petit nombre, le $+1$, et elle explique quatre étapes.

La règle $a_n = a_{n-1}^2 + c$ échoue~: $2^2 + c = 3$ impose $c = -1$, alors que $3^2 + c = 7$ impose $c = -2$. Une règle linéaire à deux termes $a_n = p\,a_{n-1} + q\,a_{n-2}$, ajustée aux premiers pas, donne 8 là où la suite donne 22.

\iqwrong{Essayer une règle de la forme $a_{n-1} \times k$ avec un $k$ croissant. Les multiplicateurs ne sont pas entiers, donc cela ne mène nulle part.}
\end{solution}

\begin{solution}{iq:iv1:sequences-patterns-and-aptitude:29}
Le nombre suivant est 132. Ce sont les nombres de Catalan\index{nombres de Catalan}, où chaque terme est le précédent multiplié par $(4n - 2)/(n + 1)$.

Les rapports valent $1$, $2$, $5/2$, $14/5$, $3$. On les réécrit sous la forme $2/2$, $6/3$, $10/4$, $14/5$, $18/6$~: les numérateurs croissent de 4 et les dénominateurs de 1. Le rapport suivant est $22/7$, et $42 \times 22/7 = 132$. La forme close est $\frac{1}{n + 1}\binom{2n}{n}$, et elle compte, par exemple, les façons d'écrire $n$ paires de parenthèses bien formées.

Les deux suites voisines échouent sur les derniers termes. Les nombres de Bell valent 1, 1, 2, 5, 15, 52, donc ils divergent à 14. La suite $(3^{n-1} + 1)/2$ donne 1, 2, 5, 14, 41 après le 1 initial, si bien qu'elle coïncide jusqu'à 14, puis donne 41 au lieu de 42.

\iqfollow{À quoi ressemble le décompte des parenthèses pour $n = 3$~?}{Cinq~: ((())), (()()), (())(), ()(()) et ()()().}
\iqwhy{C'est l'un des cas où le rappel de suites connues aide, mais ce sont les rapports qui le démontrent. Le même comptage revient dans les questions sur les expressions bien parenthésées et les permutations de pile en entretien de programmation.}
\end{solution}

\begin{solution}{iq:iv1:sequences-patterns-and-aptitude:30}
Le nombre suivant est 102.

Chaque terme est le précédent plus le produit de ses chiffres non nuls\index{suites : règles sur les chiffres}. Le début ressemble à un doublement, et c'en est un, car pour un nombre à un chiffre, le produit des chiffres est le nombre lui-même. À partir de 16~: $16 + 1 \times 6 = 22$, $22 + 2 \times 2 = 26$, $26 + 2 \times 6 = 38$, $38 + 3 \times 8 = 62$, $62 + 6 \times 2 = 74$, et $74 + 7 \times 4 = 102$.

Le doublement est un leurre. Il échoue dès 22, et les différences qui suivent, 6, 4, 12, 24, 12, n'ont de régularité qu'au niveau des chiffres. Un pas qui diminue puis augmente ainsi, dans une suite par ailleurs croissante, désigne les chiffres. Après avoir essayé la somme des chiffres, qui donne $16 + 7 = 23$ et échoue, on essaie ensuite le produit.

\iqfollow{Pourquoi la règle dit-elle «~non nuls~»~?}{Le pas suivant en fait usage~: à partir de 102, le produit de tous les chiffres vaudrait 0 et la suite stagnerait. Avec des chiffres non nuls, le pas vaut 2, donc le terme suivant est 104.}
\iqwrong{Ajuster une récurrence du type «~doubler, puis ajouter un peu~». Elle utilise un nombre libre à chaque pas, et n'est donc pas une règle du tout.}
\end{solution}
